\documentclass[12pt]{article}
\usepackage[brazilian]{babel}
\usepackage{csquotes}
\usepackage{glossaries}
\usepackage[
  perfil=artigo,
  artigo-colunas=duas,
  citacoes=autor-data
]{abntex3}

\makeglossaries
\newglossaryentry{pipeline}{
  name={pipeline},
  description={sequência automatizada de etapas de construção e validação}
}
\addbibresource{referencias-exemplo.bib}
\ABNTEXsetup{
  titulo={Modelo canônico completo de artigo em duas colunas},
  subtitulo={variante editorial auditável}
}
\ABNTEXartigoSetup{
  titulo-outro-idioma={Complete canonical two-column article template},
  resumo={Este artigo exercita a variante editorial em duas colunas e todos os elementos aplicáveis atualmente implementados.},
  palavras-chave={artigo; duas colunas; auditoria; LaTeX},
  resumo-outro-idioma={This article exercises the two-column editorial variant and every currently implemented applicable element.},
  palavras-chave-outro-idioma={article; two columns; audit; LaTeX},
  data-submissao={12 de fevereiro de 2026},
  data-aprovacao={22 de abril de 2026},
  identificacao-disponibilidade={DOI: 10.0000/exemplo.duas-colunas}
}
\ABNTEXartigoautor
  {Ana Silva}
  {Doutora em Ciência da Informação}
  {Universidade Exemplo}
  {ana@example.org}
\ABNTEXartigoautor
  {Bruno Souza}
  {Mestre em Preservação Digital}
  {Instituto Exemplo}
  {bruno@example.org}

\setlength{\emergencystretch}{1em}

\begin{document}
\ABNTEXartigopretextual
\ABNTEXartigotextual

\ABNTEXartigointroducao[Contextualização]{%
  O título adaptado demonstra que a nomenclatura textual pode ser escolhida
  pelas autorias. O estudo examina um \gls{pipeline} reprodutível
  \parencite{associacao2025}.
}
\ABNTEXartigodesenvolvimento[Procedimentos, resultados e discussão]{%
  Os procedimentos incluem inventário, compilação e inspeção.
  \begin{ABNTEXcitacaolonga}
  Este trecho permite observar recuo, corpo e espaço em uma coluna estreita.
  O texto é apenas demonstrativo e deve ser substituído pela fonte real.
  \end{ABNTEXcitacaolonga}
  \ABNTEXnota{Nota demonstrativa da variante em duas colunas.}

  \subsection{Resultados}
  \begin{figure}[ht]
    \centering
    \rule{0.8\columnwidth}{1.5cm}
    \caption{Pipeline demonstrativo}
    \ABNTEXfonte{elaborada pelas autorias.}
  \end{figure}
  \begin{table}[ht]
    \centering
    \caption{Validações realizadas}
    \begin{tabular}{ll}
      Item & Estado\\
      Colunas & duas\\
      Fluxo & completo
    \end{tabular}
    \ABNTEXfonte{elaborada pelas autorias.}
  \end{table}

  \subsection{Discussão}
  \begin{sloppypar}
  A escolha de colunas é editorial e não altera os elementos estruturais.
  A mesma sequência semântica permanece disponível nas duas variantes, o que
  permite comparar a apresentação sem confundir estrutura e projeto gráfico.
  \end{sloppypar}
}
\ABNTEXartigoconsideracoesfinais[Conclusões]{%
  Os objetivos foram atendidos e a variante permanece verificável pelos
  mesmos comandos públicos.
}

\ABNTEXartigopostextual
\ABNTEXbibliografia
\ABNTEXglossario
\ABNTEXappendix{Lista de verificação elaborada pelas autorias}
Material produzido para a auditoria do artigo.
\ABNTEXannex{Documento externo de demonstração}
Material não produzido pelas autorias.
\ABNTEXartigoagradecimentos{%
  À comunidade e às pessoas responsáveis pela avaliação independente.
}
\end{document}
